\section{Cable and Overhead-Line Parameter Uncertainty}
\label{sec:parameter_uncertainty}

% $\mathcal{P}$ and $\bar{p}$ are \ParamSet and \ParamVec of the mathematical framework; use the macros once they move to nomenclature.tex.

The per-unit-length series impedance~$\bar{z}'_{\Line}(\Frequency;\bar{p})$ and shunt admittance~$\bar{y}'_{\Line}(\Frequency;\bar{p})$ follow from physical inputs of the cables and overhead lines, accounting for the corresponding material and dimensional tolerances, and computed with the Julia package \textsc{LineCableModels.jl}~\cite{lcm_2026}. Each cable or line is a system of tubular conductors separated by insulating and semiconducting layers, in the air or in the earth~\cite{ametani_cable_2015}. For~$n$ conductors, the $n \times n$ matrices of series impedance and shunt admittance per unit length are
\begin{IEEEeqnarray}{ r C l }
    Z'(\Frequency;\bar{p}) &=& Z'_{\text{c}} + Z'_{\text{i}} + Z'_{\text{e}}, \label{eq:pul_impedance} \\
    Y'(\Frequency;\bar{p}) &=& \mathrm{j} 2\pi\Frequency \left( P'_{\text{i}} + P'_{\text{e}} \right)^{-1}, \label{eq:pul_admittance}
\end{IEEEeqnarray}
in which $Z'_{\text{c}}$ collects the surface impedances of the conductors, $Z'_{\text{i}}$ the impedances of the insulating layers and $Z'_{\text{e}}$ the earth-return impedances; and $P'_{\text{i}}$ and $P'_{\text{e}}$ are the potential coefficients of the insulating layers and of the surrounding air and earth. \textcolor{black}{These expressions hold from the fundamental to frequencies well above 2.5\,kHz, the 50th harmonic~\cite{ametani_cable_2015}. Specifically, the formulations adopted in LineCablemodels.jl have been validated against the finite element method (FEM) up to 100 MHz.}

% The surface impedances are ratios of modified Bessel functions of~$r\sqrt{\mathrm{j}2\pi\Frequency\mu/\rho}$ for a conductor of radius~$r$, permeability~$\mu$ and resistivity~$\rho$. The factor~$\sqrt{\pi\Frequency\mu/\rho}$ is the reciprocal of the skin depth, which in aluminum falls from 12.0\,mm at 50\,Hz to 3.3\,mm at 650\,Hz, below the 13.25\,mm radius of a 500\,mm$^{2}$ core. The earth return penetrates 712\,m at 50\,Hz and 197\,m at 650\,Hz into a 100\,$\Omega\cdot$m earth, so the zero-sequence parameters, whose currents return through screens, earth wires and earth, vary most. 
The effect of temperature~$\vartheta$, in $^{\circ}$C, on metal resistivity is described by 
\begin{equation}
    \rho(\vartheta) = \rho_{20} \left[ 1 + \kappa_{20} \left( \vartheta - 20 \right) \right],
    \label{eq:resistivity_temperature}
\end{equation}
where $\rho_{20}$ is the resistivity at 20\,$^{\circ}$C and $\kappa_{20}$ its temperature coefficient~\cite{iec62641_2022,iec60287_1_1_2006}. In overhead lines, temperature also sets the sag through the elastic thermal change-of-state equation, and each conductor lies at its support height less two thirds of the sag.

Cable screens are assumed bonded and solidly earthed at both cable terminals and the earth wires at every tower, with circuits ideally transposed. This yields circulant-reduced $3 \times 3$ matrices, expressed as zero-, positive- and negative-sequence parameters ~$\bar{z}'_{\Line}$ and~$\bar{y}'_{\Line}$ of Section~\ref{sec:parametric_network}.

Table~\ref{tab:parameter_box} summarizes the cables, lines and earth of the IEEE 30-bus case, whose cable sizes and ACSR codes follow from the branch currents, rated under IEEE~738~\cite{ieee738_2023}. The intervals are the acceptance limits of the standards cited in Table~\ref{tab:parameter_box}, operating and environmental ranges for the temperatures and the earth, and study intervals otherwise. Two continuous maps, equal to the identity at feasible points, keep every point of the box physical: the phase temperature is at least the ambient one, and the wire-diameter deviations of each ACSR code are scaled back onto the diameter and mass limits of EN~50182~\cite{en50182_2001}.

Physical inputs are represented as parameter space coordinates per input of a family, e.g., one conductor temperature for all 22 cables or, for a per-size quantity, it moves every size within its own interval. Of 39 candidate inputs, the 23 that change some sequence impedance or admittance by at least 0.1\,\% at \textcolor{black}{50, 650 or 1250\,Hz}, when moved alone to an end of their interval, form the \textcolor{black}{box coordinates}. The complete dataset, with the source and class of every value, and the specifications and case files are available in the HarmonicPowerModels.jl repository~\cite{hpm_2026}.
